%% OEIS Program Synthesis, Verified Semantics, and Graded Similarity
%%
%% An account of the OEIS-synthesis line of work: a verified domain-specific
%% language for integer-sequence programs, a neural synthesizer whose contracts are the
%% prover's theorems, and a graded (probabilistic-logic) similarity layer over the corpus.
%%
%% Build:
%%   pdflatex OEIS.tex

\documentclass[11pt]{article}
\usepackage[T1]{fontenc}
\usepackage[utf8]{inputenc}
\usepackage{lmodern}
\usepackage{xcolor}
\usepackage{geometry}
\geometry{margin=1in}
\usepackage{hyperref}
\usepackage{booktabs}
\usepackage{longtable}
\usepackage{listings}
\usepackage{array}
\usepackage{amsmath}
\usepackage{amssymb}
\usepackage{amsthm}
\usepackage{needspace}
\hypersetup{
  colorlinks=true,
  linkcolor=blue!55!black,
  citecolor=blue!55!black,
  urlcolor=blue!55!black
}
\setlength{\emergencystretch}{3em}

\lstset{
  basicstyle=\ttfamily\small,
  breaklines=true,
  frame=single,
  columns=fullflexible
}

\newcommand{\code}[1]{\texttt{#1}}
\newcommand{\sig}{\ensuremath{\sigma}}
\newcommand{\formalcommit}{26dd10a61a969abea0bbae023e70174491e2a946}
\newcommand{\formalblob}{https://github.com/zariuq/MeTTapedia/blob/\formalcommit/lean/mettapedia/Mettapedia}
\newcommand{\leansource}[3]{%
  \href{\formalblob/#1\#L#2}{\nolinkurl{#3}}}
\newtheorem{obligation}{Obligation}
\newtheorem{property}{Property}

\title{OEIS Program Synthesis, Verified Semantics, and Graded Similarity:\\
       A Unified Account}
\author{Oru\v{z}i \and Zarathustra Goertzel}
\date{\today}

\begin{document}
\maketitle

\begin{abstract}
We collect an ongoing line of work built around the On-Line Encyclopedia of Integer
Sequences (OEIS) as a testbed for verified program synthesis. Three strands meet here.
First, a small domain-specific language (DSL) of integer-sequence programs is given a
machine-checked operational semantics in Lean, together with a token \emph{recognizer},
a big-step generative-semantics (GSLT) layer with modal emission specifications, and a
family of formal \emph{obligations} that a synthesizer must respect. Second, a neural
sequence-to-graph synthesizer is organized so that the network's operational
contracts---a legality mask, a canonical form, relabeling invariance---have
machine-checked formal counterparts. Third, a graded similarity layer expressed in
probabilistic-logic truth values $\langle\text{strength},\text{confidence}\rangle$
recovers structure that an exact deduplication hash discards: it audits train/test
leakage, while a conservative abstract interpreter proposes totality, sign, and parity
invariants for each program. We report concrete measurements over a corpus of
$104{,}083$ synthesized programs, distinguish empirical validation of that interpreter
from kernel proof, and show how verified properties could safely guide search and
generate proof obligations.
\end{abstract}

\paragraph{Open artifacts.}
The \href{https://github.com/zariuq/MeTTapedia}{formal repository},
\href{https://github.com/zariuq/MeTTapedia/blob/main/papers/OEIS.tex}{paper source},
\href{https://github.com/zariuq/MeTTapedia/blob/main/papers/OEIS.pdf}{compiled PDF}, and
\href{https://github.com/zariuq/neuralnetworks}{neural experiment repository} are public.
The complete
\href{https://github.com/zariuq/MeTTapedia/blob/\formalcommit/lean/mettapedia/Mettapedia/GSLT/LanguageDef/Gauthier/Adjudications33.lean}{33-program adjudication module}
and the
\href{https://github.com/zariuq/MeTTapedia/tree/\formalcommit/lean/mettapedia/Mettapedia/MachineLearning/SearchGuidance/ProgramDiscovery}{provenance-aware program-discovery development}
are linked separately. The theorem links below are immutable GitHub source permalinks to
commit \code{\formalcommit}; they open the complete statement and proof, not a search page
or a moving branch location.

\section{Introduction}

The OEIS~\cite{oeis} is a natural benchmark for program synthesis: each entry is an integer sequence,
and a solution is a program that reproduces its terms. Recent self-learning systems
synthesize such programs at scale from a small combinator DSL, discovering programs for a
large fraction of the OEIS without human-written training
data~\cite{gauthier2023,gauthier2025}. Our contribution is not another synthesizer but a
\emph{verified and measured} account of the objects a synthesizer manipulates.

We separate three concerns that are often conflated:
\begin{itemize}
  \item \textbf{What a program \emph{is}} --- its syntax over a fixed operator vocabulary,
        and the machine-checked semantics of running it (Sections~\ref{sec:dsl}--\ref{sec:gslt}).
  \item \textbf{What a program \emph{does on a probe}} --- its input/output behaviour on a
        finite window (its \emph{signature} \sig), a finite observational summary
        (Section~\ref{sec:extensional}).
  \item \textbf{What properties static analysis can certify} --- totality, sign, and
        parity: intensional information inferred conservatively rather than read off
        finitely many outputs (Section~\ref{sec:intensional}).
\end{itemize}
The payoff of keeping these apart is a clean division of labour: the probe layer
gives cheap graded evidence but can never certify behaviour beyond the probe window; the
intensional layer proposes conservative all-input certificates, while the Lean layer can
turn selected claims into kernel proofs. Combined, they sharpen both benchmark hygiene and
the discovery of conjectures.

\section{The OEIS synthesis DSL}
\label{sec:dsl}

Programs are postfix token streams over a $16$-operator vocabulary evaluated on a
two-register integer state. The operators are constants \code{zero}, \code{one},
\code{two}; the index variable \code{x} $=n$ and an inner variable \code{y}; binary
arithmetic \code{addi}, \code{diff}, \code{mult}, \code{divi}, \code{modu}; a conditional
\code{cond} (test $\le 0$); bounded iterators \code{loop} and \code{loop2}; an unbounded
search \code{compr} (the $a$-th $n\ge 0$ satisfying a predicate); and stack operators
\code{push}, \code{pop}. Table~\ref{tab:ops} summarizes the arities and the partiality
character of each operator, which is decisive for the analysis of
Section~\ref{sec:intensional}.

\begin{table}[t]
\centering
\begin{tabular}{@{}llll@{}}
\toprule
operator & arity & role & partiality / termination \\
\midrule
\code{zero},\code{one},\code{two} & $0$ & constants & total \\
\code{x},\code{y} & $0$ & variables ($n\ge 0$; inner $\ge 0$) & total \\
\code{addi},\code{diff},\code{mult} & $2$ & arithmetic & total (resource: overflow) \\
\code{divi},\code{modu} & $2$ & division / modulo & \emph{genuine}: error on zero divisor \\
\code{cond} & $3$ & branch on $\le 0$ & total given branches \\
\code{loop},\code{loop2} & $3,5$ & bounded iteration & terminating (resource: timeout) \\
\code{compr} & $2$ & unbounded search & \emph{genuine}: may not terminate \\
\code{push} & $2$ & stack push & total (resource: push limit) \\
\code{pop} & $1$ & stack pop & \emph{genuine}: error on empty \\
\bottomrule
\end{tabular}
\caption{The synthesis DSL. Only \code{divi}/\code{modu} (zero divisor), \code{pop}
(empty stack), and \code{compr} (non-terminating search) introduce \emph{genuine}
partiality; \code{loop}/\code{loop2}, overflow, and timeouts are \emph{resource} limits.
This three-way character drives the totality analysis.}
\label{tab:ops}
\end{table}

A closed form illustrates the language: the factorial $n!$ is
\code{loop(mult(x,y),\,x,\,one)} --- fold multiplication by the running counter $n$ times
starting from $1$. The counting sequence $a(n)=n+1$ is \code{addi(x,one)}.

\section{Verified GSLT semantics}
\label{sec:gslt}

The DSL is given a machine-checked operational semantics in Lean. On top of it we build a
generative-semantics layer as a triple $(T,E,R)$: terms $T$, an equational theory $E$
(a setoid), and a rewrite relation $R$ with left/right coherence laws. The coarse
\emph{big-step GSLT} rewrites a program to the integer value its fuel-bounded evaluator
emits; a completed-value term is never reflexively a program, so non-triviality is
load-bearing. Two modalities express synthesis goals:
\begin{align*}
  \Diamond(\text{emits } v \text{ at } k) &\iff \exists\, s',\ \code{eval}\ \mathit{fuel}\ \mathit{sig}\ p\ (\code{seed}\ k)\ = \code{some}(v,s'),\\
  \code{EmitsPrefix}\ p\ \ell &\iff \forall i\, v,\ \ell_i = v \Rightarrow \Diamond(\text{emits } v \text{ at } i).
\end{align*}
\code{EmitsPrefix} is precisely the OEIS synthesis specification: ``$p$ reproduces the
observed prefix $\ell$.''

\paragraph{Obligations.} A synthesizer that claims correctness must respect a family of
formal obligations. We single out three that are discharged as machine-checked theorems,
axiom-clean (each theorem's \code{\#print axioms} footprint lies within
$\{\text{propext}, \text{Classical.choice}, \text{Quot.sound}\}$):
\begin{obligation}[Recognizer completeness / mask equivalence]
The legality predicate \code{maskLegal} on token streams coincides with the executable
recognizer (\leansource{GSLT/LanguageDef/Gauthier/E1.lean}{406-L415}{recognize_iff_maskLegal}),
and every well-formed program is recognized
(\leansource{GSLT/LanguageDef/Gauthier/E1.lean}{492-L499}{recognize_complete}).
Hence a constrained decoder that emits only mask-legal tokens
can generate exactly the well-formed programs.
\end{obligation}
\begin{obligation}[Relabeling invariance]
Evaluation is invariant under a signature isomorphism
(\leansource{GSLT/LanguageDef/Gauthier/E1.lean}{1612-L1628}{eval_relabel_signature_iso}):
renaming primitives consistently does not change the
computed sequence. Behaviour depends on structure, not on operator names.
\end{obligation}
\begin{obligation}[Canonicalization soundness]
Flag-driven canonicalization preserves semantics; the relevant proofs are
\leansource{GSLT/LanguageDef/Gauthier/E1.lean}{1431-L1437}{canonicalize_sound} and
\leansource{GSLT/LanguageDef/Gauthier/E1.lean}{1816-L1823}{orgCanonicalize_sound}.
Thus, for example, sorting the
children of commutative operators cannot make canonical-form deduplication merge
behaviourally distinct programs.
\end{obligation}
Further obligations include equational-theory soundness and evaluation congruence for the
modalities. The machine-checked
\leansource{GSLT/LanguageDef/Gauthier/ProbeRigidity.lean}{73-L79}{finite_probe_rigidity_negative_witness}
already settles finite-probe rigidity in the negative: two programs can agree on a finite
probe while differing extensionally. This is
the formal shadow of the leakage phenomenon in Section~\ref{sec:extensional}. A separate
small-step account remains future work; the present rewrite layer deliberately uses the
evaluator as its big-step relation.

\section{The synthesizer as specification}
\label{sec:synth}

The neural synthesizer maps a sequence prefix to a program graph (sequence-to-graph). Its
decoder is \emph{execution-guided}: at each step a legality mask restricts the token set to
those that can extend a well-formed program, and the mask is exactly the recognizer of
Obligation~1. Successive training variants add retrieval and contrastive objectives that
consult a similarity index. The
central design stance is that the network's operational contracts are the prover's
theorems: the mask's soundness, the canonical form's soundness, and relabeling invariance
are not engineering conveniences but the denotational safety envelope of the synthesizer.

\section{Leakage-safe data and the probe-behaviour layer}
\label{sec:extensional}

\paragraph{Exact deduplication and splits.} Each program's probe signature \sig{} is its
output vector on $n=0,\dots,31$ at a fixed fuel, with $\bot$ (\code{None}) marking a
non-value result at that budget. Programs with identical \sig{} form an exact
\emph{probe-signature} class (called an ``E-class'' in the implementation); this is not
extensional program equivalence. The train/val/test split is assigned by whole connected
component so that no exact-signature duplicate straddles the split boundary. Over
$104{,}083$ programs this gate is leak-tight for the measured signatures: all $3{,}007$
exact-\sig{} duplicate pairs lie in the same split.

\paragraph{Graded similarity.} Exact equality is brittle. We express graded similarity as a
probabilistic-logic truth value $\langle s,c\rangle$~\cite{pln2009}: the probe-level
strength is behavioural agreement over the \emph{defined} domain,
$s = |{\text{agree}}| / |{\text{informative}}|$, and confidence follows the chosen
evidence-count law $c = m/(m+k)$, with $m$ the number of informative probe positions and a
fixed evidence scale $k>0$. This is an evidence-weighting convention, not an empirical
calibration claim. For fully defined signatures, exact equality gives $s=1,\ m=32$; when
$\bot$ occurs, $m$ counts only informative positions. A locality-sensitive index over
signatures generates near-equivalent candidate pairs without an $O(n^2)$ scan.

\paragraph{Soft-leakage audit.} The graded layer surfaces \emph{near}-duplicates the exact
gate cannot see. Table~\ref{tab:leak} reports pairs whose two programs sit in different
splits. The benchmark-integrity figure is the fraction of evaluation programs with a near
behavioural twin in \emph{train}: at strength $\ge 0.95$, $1.84\%$ of test and $1.84\%$ of
val programs (lower bounds---the index can miss a pair, never invent one). A
disagreement-locus analysis separates genuine near-twins (disagreements confined to a
single term, none in the tail) from ``false friends'' that agree only on a shared trivial
tail; the trustworthy signal concentrates in the $\ge 0.95$ band. This vindicates gating on
\emph{exact} full-signature equality rather than a similarity threshold, which would either
over-merge on trivial tails or miss single-term twins.

\begin{table}[t]
\centering
\begin{tabular}{@{}lrr@{}}
\toprule
strength band & cross-split pairs & train$\leftrightarrow$test pairs \\
\midrule
$0.80$--$0.90$ & $17{,}250$ & $6{,}365$ \\
$0.90$--$0.95$ & $2{,}614$  & $909$ \\
$0.95$--$0.99$ & $502$      & $150$ \\
\bottomrule
\end{tabular}
\caption{Soft leakage: near-equivalent program pairs ($d\ge 1$, not byte-identical) placed
in different splits by the exact gate. Distinct evaluation programs with a train near-twin:
$1.84\%$ (test and val) at $\ge 0.95$; $5.4$--$5.8\%$ at $\ge 0.90$.}
\label{tab:leak}
\end{table}

\section{The intensional layer: static properties and their proof boundary}
\label{sec:intensional}

Where the probe layer reads behaviour off finitely many outputs, the intensional layer
uses abstract interpretation~\cite{cousot1977} to infer candidate all-input invariants.
Its transfer functions are designed to over-approximate the concrete operators, and loops
are handled by Kleene fixpoints in finite lattices. The current analyzer is implemented in
Python and tested against the evaluator; its transfer functions have not yet been proved
sound against the Lean semantics. We therefore distinguish its conservative static
certificates from the kernel-checked program-correctness proofs of
Section~\ref{sec:adjudication}. We compute three properties.

\begin{property}[Totality]
A program is \emph{classified total} when only resource limits can stop it: it contains
neither \code{compr} nor \code{pop}, and every \code{divi}/\code{modu} divisor is certified
nonzero. This is a conservative lower bound---``unknown'' does not mean partial.
\end{property}
\begin{property}[Sign]
A sign-domain analysis (with $\code{x}=n\ge 0$, and $a\bmod b\ge 0$ for $b>0$) certifies
the OEIS keyword \code{nonn} (nonnegativity) for many programs under its abstract rules.
\end{property}
\begin{property}[Parity]
A $\mathbb{Z}/2$ analysis certifies ``always even'' / ``always odd'' under its abstract
rules (e.g.\ \code{mult} by an even factor is even).
\end{property}

\paragraph{Validation gates.} Two independent finite checks test the implementation; they
do not replace a soundness proof. A static
gate compares every sign/parity claim against the real signatures: $0$ violations across
all $104{,}083$ programs ($\approx 3.3$M value-checks). An empirical totality gate re-runs
the real evaluator on $500$ programs classified total to $n=48$ (beyond the probe window):
$0$ genuine errors, only resource timeouts. Table~\ref{tab:props} gives the distribution.

\begin{table}[t]
\centering
\begin{tabular}{@{}lrr@{}}
\toprule
static property & count & share \\
\midrule
total          & $37{,}996$ & $36.5\%$ \\
nonnegative (\code{nonn}) & $45{,}251$ & $43.5\%$ \\
positive       & $14{,}907$ & $14.3\%$ \\
even           & $2{,}725$  & $2.6\%$ \\
odd            & $2{,}926$  & $2.8\%$ \\
\bottomrule
\end{tabular}
\caption{Static-property distribution over $104{,}083$ programs. Shares are lower bounds
under the analyzer's intended conservative interpretation.}
\label{tab:props}
\end{table}

\paragraph{Disambiguating the signature.} A $\bot$ in \sig{} is ambiguous: is the program
genuinely partial, or merely slow at the probe budget? Under the analyzer's intended
soundness, a total classification attributes it to resource exhaustion. Of the $12{,}966$
programs carrying a $\bot$, $2{,}345$ are classified total, while $10{,}621$ are not, with
attributed cause: $6{,}539$ \code{compr} (non-terminating search), $3{,}023$ \code{pop}
(empty list), $1{,}059$ \code{divi}/\code{modu} (possible zero divisor). The intensional
layer thereby partitions a signal that finite execution alone leaves ambiguous.

\section{Filtering and measuring in the OEIS loop}
\label{sec:filter}

Both roles suggested by Table~\ref{tab:props} are useful inside a synthesis loop.

\paragraph{Filtering.} An OEIS target often carries stated properties---for example its
keywords (\code{nonn}, \code{sign}, \code{mult}) or invariants justified by its defining
formula. A candidate can be pruned soundly only when two premises are established: the
target invariant is itself justified, and the analyzer proves an incompatible candidate
invariant. An ``unknown'' abstract result must remain legal; requiring every candidate to
be \emph{proved} nonnegative would be incomplete and could discard a correct program.
Thus, before the analyzer is connected to the decoder, uncertain properties may rank
candidates, while hard pruning must be restricted to proved contradictions. A future Lean
soundness theorem for the analyzer would make that restricted semantic mask
recall-preserving on its stated domain.

\paragraph{Measuring.} The same properties are theoretical measurements of a program and of
a sequence. Totality analysis distinguishes a possibly partial sequence from one classified
as slow-to-compute but total; the sign/parity results are static-analysis counterparts of
OEIS keywords, so annotations can be predicted and checked at scale. More finely, the
graded truth values give each program pair an evidence-weighted similarity, and formal
proofs can supply factive intensional features. Revision is licensed only for
source-disjoint packets: a proof must not be counted again as independent support for the
translation or observations on which it depends. The corpus thereby becomes not merely
programs and outputs, but programs, outputs, provenance, and proofs.

\section{The loop compounds: a three-seed cumulative replication}
\label{sec:loop}

The construction above is only interesting if the loop it serves actually runs. It does. We
report a prospectively registered replication of the self-learning regime
of~\cite{gauthier2023} using a typed-action decoder over the verified skeleton of
Section~\ref{sec:synth}: the network never emits program text, only \emph{legal construction
actions} over an exact state, and a deterministic checker decides acceptance.

The protocol is cumulative. Each of three independently seeded worlds---with random seeds
controlling initialisation, minibatch order, dropout, and search
tie-breaking; the data pool ($79{,}249$ rows) and the $240$ search tasks are shared---trains
on its accumulated verified corpus, searches under a fixed budget, checks every candidate
against the whole OEIS, adds newly covered sequences to its corpus, and repeats. The
endpoint is cumulative coverage $C_g$ and its increment $D_g = C_g - C_{g-1}$; no
comparison against held-back external data governs continuation, since in the deployed
regime no such reservoir exists.

The shared initial corpus is not chosen freely: it is the reference system's own public
bootstrap seed \code{sol0}~\cite{gauthier2023}, $3{,}771$ sequences found by \emph{random}
program search with no learning, and the very origin from which that system subsequently
grew its full published corpus through many iterations of the loop we are running here.
Starting from that seed rather than from a mature corpus is deliberate, with three
consequences. The origin is reproducible and learning-free, a fixed public artefact anyone
can replay from. Our run is then the \emph{same} self-learning process re-executed---not a
fork of another system's answers---so a rediscovery is a genuine held-out generalisation
within our own lineage rather than a lookup. And the base stays small enough that a world's
own discoveries form a measurable share of the training signal instead of being drowned by
it, which is what makes the loop's dynamics observable at all. Coverage per world:

\begin{center}
\begin{tabular}{lrrr}
\toprule
world & initial & after gen.\ 1 & after gen.\ 2 \\
\midrule
17 & 3{,}771 & 3{,}789 & 3{,}846 \\
29 & 3{,}771 & 3{,}841 & 3{,}891 \\
43 & 3{,}771 & 3{,}807 & 3{,}862 \\
\bottomrule
\end{tabular}
\end{center}

At generation three we registered, before any update, a three-way comparison of how the
accumulated discoveries enter training: \emph{uniform} replay (literal pooling, under which
the $194$ cumulative discoveries receive their natural $0.24\%$ of sampling mass);
\emph{balanced} replay ($25\%$ discovery mass); and balanced replay with the auxiliary
predictive heads disabled. All three arms continue from the same per-world parent
checkpoint under identical update counts and search budgets, and are judged solely by new
whole-OEIS A-numbers beyond the pooled $3{,}986$ already known. New A-numbers, generation
three:

\begin{center}
\begin{tabular}{lrrrr}
\toprule
setting & world 17 & world 29 & world 43 & pooled \\
\midrule
uniform replay & 48 & 39 & 72 & 159 \\
balanced replay, aux.\ off & 65 & 155 & 88 & 308 \\
balanced replay, aux.\ on & 135 & 155 & 177 & \textbf{467} \\
\bottomrule
\end{tabular}
\end{center}

Every setting is positive in every world: nine of nine cells extend coverage, so the loop
compounds rather than saturating at this scale. Two registered contrasts separate the
mechanism. Balanced replay is worth $+308$ pooled over uniform---at literal pooling the
model's own discoveries are numerically drowned by the base corpus, and the loop is starved
of exactly the signal it exists to generate. The auxiliary heads are worth a further $+159$.
The generation-three increment under the selected setting exceeds the generation-two
increment by roughly threefold, and the setting carried into generation four was fixed by
the registered selection rule rather than chosen after inspection.

Two negative results bound the claim, and we record them because they materially limit its
interpretation. First, an ablation asking whether a world's own discoveries beat an equal mass
of \emph{unused external certified examples} answers no, in all three worlds
($162$ versus $253$ pooled): feedback submitted more candidates yet explored fewer distinct
programs ($5{,}540$ versus $6{,}233$), a diversity contraction rather than a fitting
failure. That comparison is a legitimate ablation and a poor objective---the fresh arm is an
oracle-advantaged comparator that does not exist in the deployed loop, where the corpus is
whatever the loop has verified. Second, no arm in any generation solved its own assigned
held-out target within budget. The evidence therefore concerns learned \emph{discovery}
guidance over the catalogue---the regime~\cite{gauthier2023} operates in---and not
target-conditioned synthesis, which remains open.

What the verified layer contributes here is attributability rather than yield. Because the
legality mask is machine-checked, a discovery cannot be an artefact of a malformed
candidate slipping past a hand-written filter; because ranking is provably
recall-preserving, guidance changes search order and cannot manufacture or destroy
solutions; and because every accepted program is checked against the whole corpus, the
increments above are exact catalogue-coverage counts under that checker rather than
estimates or full-domain correctness claims. The loop's growth and the loop's
trustworthiness are established by different machinery, and both are needed.

\subsection{The compounding claim, tested against its own control}
\label{sec:loopcontrol}

The replication above shares a weakness with the regime it replicates: cumulative coverage
rises, but training time rises with it, so growth and continued optimisation are
confounded. We therefore ran the compounding question separately, on a faithful port of
the reference system's own architecture---a $10{,}781{,}697$-parameter scaled-Luong
bidirectional-LSTM sequence-to-sequence model---under the reference system's own protocol,
in which every generation is trained \emph{from fresh initialisation} and knowledge is
carried by the accumulated corpus rather than by the weights.

That protocol makes an exact control available. In two independently initialised worlds,
one arm trains generation two on the seed corpus \emph{augmented} with generation one's own
verified discoveries; the paired control trains on the seed corpus \emph{alone} with fresh
initialisation. Both are compute-matched, decoded and whole-OEIS-checked identically, and
scored against the same fixed baseline of $5{,}631$ known sequences, so the sole difference
between them is whether the model's own discoveries entered its training data. That
difference is the loop.

\begin{center}
\begin{tabular}{lrr}
\toprule
world & discoveries fed back & seed corpus only \\
\midrule
1 & $1{,}980$ & $375$ \\
2 & $2{,}051$ & $519$ \\
\bottomrule
\end{tabular}
\end{center}

Feeding verified discoveries back finds roughly four to five times the new frontier that
re-seeding on the same corpus finds, in both worlds, with the pre-registered ordering met
in each. Every count carries a content-addressed job identity and was reproduced by an
independent verification pass. Because both arms start from fresh weights, no part of this
gap can be attributed to accumulated optimisation: the corpus alone carries it.

The contrast with the replication above is itself informative, and we state it plainly
rather than let it pass as an unnoticed deviation. The reference implementation
re-initialises its network every generation; the typed-action campaign of this section
continues weights and optimiser state across generations. The two designs answer different
questions---one isolates the corpus, the other measures a lineage---and only the first can
attribute compounding to feedback without qualification.

\subsection{Eight generations, two worlds, four state carriers}
\label{sec:eightgen}

The lineage design was then run to its registered endpoint: two independently initialised
worlds, eight generations each, four matched-capacity decoder state carriers over the
shared typed-action interface---a gated recurrent vector, a typed-slot workspace, and two
belief-register variants (retention derived from the drift statistics versus learned).
Every candidate was checked against the whole encyclopedia; all $64$ cells passed
independent verification. Distinct new sequences across everything: $984$, from
near-identical world totals ($639$ and $638$, sharing only $293$).

\begin{center}
\begin{tabular}{lrrr}
\toprule
state carrier & distinct & found by no other & training cost \\
\midrule
recurrent vector (GRU) & 435 & 241 & $1.00\times$ \\
typed workspace        & 383 & 171 & $2.85\times$ \\
belief, derived retention & 294 & 119 & $7.94\times$ \\
belief, learned retention & 244 & 138 & $3.53\times$ \\
\bottomrule
\end{tabular}
\end{center}

Three readings, in decreasing order of certainty. First, the yield ordering tracks the
\emph{decisiveness} ordering exactly---the recurrent vector had the lowest action entropy,
the highest gold-action probability, and the fewest duplicate proposals---which is the
mechanism the
\leansource{MachineLearning/SearchGuidance/DecisivenessYield.lean}{242-L248}{equal_crossEntropy_strictly_ordered_expectedYield}
theorem isolates under a fixed search budget. The broader fixed-budget scope and its
counterexamples are collected by
\leansource{MachineLearning/SearchGuidance/BudgetedCoverageCrown.lean}{102-L130}{budgetedSearchYield_crown}.
Second, $669$ of $984$ sequences were found by exactly one state carrier, with pairwise
overlaps of $8$--$16\%$: the carriers steer search into substantially different territory,
so the union is worth far more than the best single arm, yet a fixed-split portfolio
under one budget reached only $399$ against the best arm's $435$---complementarity is
real and naive allocation fails to harvest it. Third, and most tentatively, the
derived-retention belief arm (retaining $\approx 94\%$ of prior evidence per revision)
out-yielded the learned-retention arm, which learned to discard roughly half its evidence
per revision---the direction the drift theory predicts in a near-stationary domain, where
the optimal retention approaches one. Margins varied across the two worlds (one a clear
recurrent-vector win, the other a near-tie), so these are descriptive statements about
two runs, not population claims. Roughly three quarters of every arm's search budget was
spent re-proposing programs it had already generated, which bounds how much any
architectural conclusion can matter before search-side deduplication is addressed.

\section{From coverage to proof: adjudicating discoveries against their definitions}
\label{sec:adjudication}

Everything above establishes that a discovered program reproduces a sequence's stored
terms. That is weaker than it sounds. The published entries are short---median $38$ terms
corpus-wide, and only $26$ for the sequences our own loop discovered, with $22\%$ resting
on fewer than twenty terms and one on three. This gap is consistent with a simple
selection mechanism: short entries impose fewer observed constraints and are therefore
more vulnerable to coincidental matches. And finite agreement is exactly what our
finite-probe rigidity
counterexample proves insufficient---agreement on a finite prefix need not force
extensional equality.

So we asked the question the loop cannot answer by itself: does the discovered program
actually compute the sequence its entry \emph{defines}? This is answerable, because the
encyclopedia carries far more than terms. Of our $984$ discoveries, all $984$ have
keywords, $88\%$ an explicit formula field, $88\%$ a reference program, and $75\%$ prose
commentary. These fields provide candidate sources for specifications, but their precision
varies: a closed formula is much less ambiguous than a keyword or informal comment.

\paragraph{Protocol.} A definition can always be made to match a program after the fact,
so the order of work is the whole methodology. Candidate programs were frozen first, from
the completed campaign, as exact generated token streams with their hashes. Sequences were
then selected on \emph{metadata alone}, and each specification was formalized and locked
\emph{before} its candidate program was revealed. The specification therefore states a
proof obligation rather than describing a known solution. This is the ordinary discipline
of verification---a specification confronting a pre-existing implementation---and not, as
one might worry, a translation from definition to program.

\paragraph{What is proved.} For a specification with domain $D$, index map, and value
function, and a frozen token stream $t$:
\[
\mathrm{Realizes}(\mathrm{spec}, p) \;\equiv\;
\forall\, i,\; D(\mathrm{index}(i)) \;\rightarrow\;
\mathrm{EventuallyEmits}\bigl(p,\, i,\, \mathrm{value}(\mathrm{index}(i))\bigr),
\]
where $\mathrm{EventuallyEmits}$ requires a fuel threshold beyond which the interpreter
\emph{stably} returns that value. Three things deserve emphasis. The quantifier is over
the entire declared domain, not a probed window, so this is extensional correctness rather
than extended testing. Termination is part of the claim: fuel stability is proved, not
assumed, and a companion theorem establishes that a successful evaluation is unique
across sufficient fuel
(\leansource{GSLT/LanguageDef/Gauthier/OEISSequenceSemantics.lean}{61-L71}{eventuallyEmits_unique}).
The formal definitions are
\leansource{GSLT/LanguageDef/Gauthier/OEISSequenceSemantics.lean}{42-L44}{EventuallyEmits},
\leansource{GSLT/LanguageDef/Gauthier/OEISSequenceSemantics.lean}{74-L76}{Realizes}, and
\leansource{GSLT/LanguageDef/Gauthier/OEISSequenceSemantics.lean}{78-L79}{CandidateRealizes}.
Finally, the verified program is not a manual Lean
reimplementation: the frozen token stream is parsed by the same recognizer the synthesizer
uses, with the parse discharged by kernel computation, so what is verified is the artifact
the network emitted.

\paragraph{Result.} Thirty-three sequences were adjudicated, and every frozen program
associated with them was proved extensionally correct on its full domain. Each carries an
axiom audit; all depend only on the standard logical basis. Each record pins the OEIS
identifier, entry revision, entry digest, offset, token stream and program hash, so a
reader can reconstruct exactly which published text was translated. The registry cardinality
is itself checked by
\leansource{MachineLearning/SearchGuidance/ProgramDiscovery/OEISAdjudicationEvidence.lean}{148-L149}{certifiedProofPacket_count},
and all 33 individual proofs are in the linked adjudication module above.

The OEIS definitions and metadata quoted below are attributed to the OEIS Foundation
Inc.~\cite{oeis}, from the pinned \code{oeisdata} snapshot at revision
\code{a6e0f22854cc1c307da428e9d6295093781df7fa}, and are used under the Creative Commons
Attribution-ShareAlike 4.0 license. Each formalization additionally stores the exact entry
digest rather than relying on the repository revision alone.

The five examples below are deliberately not compressed into one table. For readability,
write $\operatorname{Loop}(f,k,z)=f^{\circ k}(z)$ for these loop bodies, all of which ignore
the loop counter, and abbreviate the Lean namespace \code{GauthierE2.P} as \code{P}. The
displayed ``semantic normal form'' is the value established by the proof, not merely a
pretty-printing or a finite simplifier run. Each separate Lean listing then shows the source
specification, frozen hash and tokens, parsed program AST, recognizer equality, literal
theorem type, and the theorem with all three semantic definitions unfolded. All five
specifications have offset zero, so the visible domain premise
$0\leq\mathrm{Int.ofNat}(\mathrm{position})$ is automatic; we retain it to expose the
domain-aware theorem exactly.

\Needspace{29\baselineskip}
\paragraph{A016779: one cubing step.}
The discovered GSLT program has the readable semantic normal form
\[
p_{28}(n)=\operatorname{Loop}(x\mapsto x^3,1,1+3n)=(3n+1)^3.
\]
\noindent Full machine-checked statement and proof:
\leansource{GSLT/LanguageDef/Gauthier/Adjudications33.lean}{248-L263}{A016779_candidate28_correct}.
\begin{lstlisting}[caption={A016779, discovered exclusively by the recurrent-vector arm
in world 1, generation 8. The proof ranges over the entire domain, not only the 36
published terms.}]
-- OEIS %N: a(n) = (3*n + 1)^3.  %O 0,2
def A016779.source := sourceOf "A016779"
  "56bc1bf16ed408a98b0d997bc8006d1b34e8796dcaf47bfbd91586155f526b97" 0
def A016779.spec := specOf 0 (fun n => (3 * n + 1) ^ 3)

def candidate28 : FrozenCandidate where
  programSha256 := "dba931cee529c62f5c8284b78fbb0974327221981a029c00f99a6b4d0711eb2d"
  tokens := [10, 10, 5, 10, 5, 1, 1, 10, 10, 3, 10, 3, 3, 9]
  program := P.loop (P.mult (P.mult P.X P.X) P.X) P.o
    (P.addi P.o (P.addi (P.addi P.X P.X) P.X))
  recognized := rfl

#check A016779_candidate28_correct :
  CandidateRealizes A016779.spec candidate28

-- The same theorem, fully unfolded and kernel-checked:
example : forall position : Nat, (0 : Int) <= Int.ofNat position ->
    exists threshold, forall fuel, threshold <= fuel ->
      GauthierE2.term fuel orgMemoSignature candidate28.program
        (Int.ofNat position) =
          some ((3 * Int.ofNat position + 1) ^ 3) := by
  simpa [CandidateRealizes, Realizes, EventuallyEmits,
    A016779.spec, specOf, SequenceSpec.index]
    using A016779_candidate28_correct
\end{lstlisting}

\Needspace{29\baselineskip}
\paragraph{A016780: two squaring steps.}
Here the network expresses a fourth power as two iterations of squaring:
\[
 p_{29}(n)=\operatorname{Loop}(x\mapsto x^2,2,1+3n)
          =\bigl((3n+1)^2\bigr)^2=(3n+1)^4.
\]
\noindent Full machine-checked statement and proof:
\leansource{GSLT/LanguageDef/Gauthier/Adjudications33.lean}{265-L280}{A016780_candidate29_correct}.
\begin{lstlisting}[caption={A016780, discovered exclusively by the belief-register arm
through derived retention in world 1, generation 4. The OEIS entry contains 31 published
terms; the theorem is universal.}]
-- OEIS %N: a(n) = (3*n + 1)^4.  %O 0,2
def A016780.source := sourceOf "A016780"
  "86c883b0debc37c8ada7ab8f7e0414d0f6f156b2533f6090d993718f45737168" 0
def A016780.spec := specOf 0 (fun n => (3 * n + 1) ^ 4)

def candidate29 : FrozenCandidate where
  programSha256 := "c57882d00078998c452c5feee21dc0c7d0399e0f60379d118fd421c54acf6d67"
  tokens := [10, 10, 5, 2, 1, 1, 2, 3, 10, 5, 3, 9]
  program := P.loop (P.mult P.X P.X) P.tw
    (P.addi P.o (P.mult (P.addi P.o P.tw) P.X))
  recognized := rfl

#check A016780_candidate29_correct :
  CandidateRealizes A016780.spec candidate29

-- The same theorem, fully unfolded and kernel-checked:
example : forall position : Nat, (0 : Int) <= Int.ofNat position ->
    exists threshold, forall fuel, threshold <= fuel ->
      GauthierE2.term fuel orgMemoSignature candidate29.program
        (Int.ofNat position) =
          some ((3 * Int.ofNat position + 1) ^ 4) := by
  simpa [CandidateRealizes, Realizes, EventuallyEmits,
    A016780.spec, specOf, SequenceSpec.index]
    using A016780_candidate29_correct
\end{lstlisting}

\Needspace{29\baselineskip}
\paragraph{A016814: one squaring step.}
The initial expression changes, while the learned loop skeleton is reused:
\[
p_{30}(n)=\operatorname{Loop}(x\mapsto x^2,1,1+4n)=(4n+1)^2.
\]
\noindent Full machine-checked statement and proof:
\leansource{GSLT/LanguageDef/Gauthier/Adjudications33.lean}{282-L297}{A016814_candidate30_correct}.
\begin{lstlisting}[caption={A016814, discovered exclusively by the belief-register arm
through derived retention in world 2, generation 1. The proof replaces a 44-term match
with a full-domain theorem.}]
-- OEIS %N: a(n) = (4*n + 1)^2.  %O 0,2
def A016814.source := sourceOf "A016814"
  "2bb3a661541a70b58c6a0e2fd6b96931502217d3eead53dd16a0d04607b4a3a2" 0
def A016814.spec := specOf 0 (fun n => (4 * n + 1) ^ 2)

def candidate30 : FrozenCandidate where
  programSha256 := "49e6cea5ab07ccf1dfac616e36e42e5e49c7ed074c227c30ddf946f24587efef"
  tokens := [10, 10, 5, 1, 1, 2, 2, 3, 10, 5, 3, 9]
  program := P.loop (P.mult P.X P.X) P.o
    (P.addi P.o (P.mult (P.addi P.tw P.tw) P.X))
  recognized := rfl

#check A016814_candidate30_correct :
  CandidateRealizes A016814.spec candidate30

-- The same theorem, fully unfolded and kernel-checked:
example : forall position : Nat, (0 : Int) <= Int.ofNat position ->
    exists threshold, forall fuel, threshold <= fuel ->
      GauthierE2.term fuel orgMemoSignature candidate30.program
        (Int.ofNat position) =
          some ((4 * Int.ofNat position + 1) ^ 2) := by
  simpa [CandidateRealizes, Realizes, EventuallyEmits,
    A016814.spec, specOf, SequenceSpec.index]
    using A016814_candidate30_correct
\end{lstlisting}

\Needspace{29\baselineskip}
\paragraph{A070439: a computed modulus.}
The program does not contain the literal $16$; it constructs it by squaring twice from two:
\[
 p_{45}(n)=n^2\bmod\operatorname{Loop}(x\mapsto x^2,2,2)
          =n^2\bmod16.
\]
\noindent Full machine-checked statement and proof:
\leansource{GSLT/LanguageDef/Gauthier/Adjudications33.lean}{175-L195}{A070439_candidate45_correct}.
\begin{lstlisting}[caption={A070439, discovered exclusively by the recurrent-vector arm
in world 1, generation 5. The synthesized denominator computes 16 rather than storing it.}]
-- OEIS %N: a(n) = n^2 mod 16.  %O 0,3  (101 published terms)
def A070439.source := sourceOf "A070439"
  "df5acb01e6c2570f6b4592782d285292d0cb69a4cd1987052c2f5ac6b0cc6de4" 0
def A070439.spec := specOf 0 (fun n => n ^ 2 % 16)

def candidate45 : FrozenCandidate where
  programSha256 := "315061126a0a0fb9eb5c6b0de9be9daa0977d662244d07020ff5276021a5f81e"
  tokens := [10, 10, 5, 10, 10, 5, 2, 2, 9, 7]
  program := P.modu (P.mult P.X P.X)
    (P.loop (P.mult P.X P.X) P.tw P.tw)
  recognized := rfl

#check A070439_candidate45_correct :
  CandidateRealizes A070439.spec candidate45

-- The same theorem, fully unfolded and kernel-checked:
example : forall position : Nat, (0 : Int) <= Int.ofNat position ->
    exists threshold, forall fuel, threshold <= fuel ->
      GauthierE2.term fuel orgMemoSignature candidate45.program
        (Int.ofNat position) =
          some (Int.ofNat position ^ 2 % 16) := by
  simpa [CandidateRealizes, Realizes, EventuallyEmits,
    A070439.spec, specOf, SequenceSpec.index]
    using A070439_candidate45_correct
\end{lstlisting}

\Needspace{29\baselineskip}
\paragraph{A070478: the cubed companion.}
The same learned construction of $16$ is paired with a cubed numerator:
\[
 p_{46}(n)=n^3\bmod\operatorname{Loop}(x\mapsto x^2,2,2)
          =n^3\bmod16.
\]
\noindent Full machine-checked statement and proof:
\leansource{GSLT/LanguageDef/Gauthier/Adjudications33.lean}{197-L221}{A070478_candidate46_correct}.
\begin{lstlisting}[caption={A070478, discovered exclusively by the recurrent-vector arm
in world 1, generation 2. Its 96 published terms become a universal stable-evaluation
claim.}]
-- OEIS %N: a(n) = n^3 mod 16.  %O 0,3
def A070478.source := sourceOf "A070478"
  "cd9f8752f701420ceb43a72e77db95d1d3816fddc6db55c8b531aabbab41947b" 0
def A070478.spec := specOf 0 (fun n => n ^ 3 % 16)

def candidate46 : FrozenCandidate where
  programSha256 := "310209b1db332a56d73abbd2476b62292b37548b20f7114124632ba8cdbb61cc"
  tokens := [10, 10, 5, 10, 5, 10, 10, 5, 2, 2, 9, 7]
  program := P.modu (P.mult (P.mult P.X P.X) P.X)
    (P.loop (P.mult P.X P.X) P.tw P.tw)
  recognized := rfl

#check A070478_candidate46_correct :
  CandidateRealizes A070478.spec candidate46

-- The same theorem, fully unfolded and kernel-checked:
example : forall position : Nat, (0 : Int) <= Int.ofNat position ->
    exists threshold, forall fuel, threshold <= fuel ->
      GauthierE2.term fuel orgMemoSignature candidate46.program
        (Int.ofNat position) =
          some (Int.ofNat position ^ 3 % 16) := by
  simpa [CandidateRealizes, Realizes, EventuallyEmits,
    A070478.spec, specOf, SequenceSpec.index]
    using A070478_candidate46_correct
\end{lstlisting}

Two details in these listings repay attention. The programs need not mirror the definitions
syntactically: A070439's specification is $n^{2} \bmod 16$, while the discovered program
computes the modulus as \code{loop (mult X X) tw tw}---sixteen obtained by iterated
squaring from two, rather than written down---and the proof must reconcile the two.
And the discovering arms are not the ones the headline count favours: two of these five
came from the belief-register architecture, which finished \emph{last} on raw yield
($294$ against the recurrent vector's $435$). All five were found by exactly one
architecture. The diversity argument of \S\ref{sec:eightgen} was until now a statement
about counts; here it is a statement about programs that are provably correct.

\paragraph{What this does and does not establish.} The cohort was selected for
formalizability, so $33/33$ is emphatically not an estimate of the correctness rate among
all $984$ discoveries; the sequences resting on three or four published terms are exactly
those excluded by that selection, and they remain the open question. More subtly, the
proofs certify that each program equals \emph{our formalization} of a definition. Whether
that formalization faithfully renders the encyclopedia's source text is a separate judgment,
and it is not one a kernel can discharge. We therefore keep two links distinct: program
$\equiv$ specification is factive; specification $\equiv$ intended sequence is a reviewed
translation carrying its own provenance. Because the proof depends on the translation, the
two must not be pooled as independent support---an independent semantic review is a
genuinely separate source, and only that can move the second link. This dependence boundary
is formalized by
\leansource{MachineLearning/SearchGuidance/ProgramDiscovery/OEISAdjudicationEvidence.lean}{112-L119}{translation_proof_have_no_additiveRevisionLicense}.

\paragraph{Refutation is generative.} A failed adjudication is not a discarded result. A
formalized definition is a certified generator: it produces additional terms beyond the
published window, with proof, turning a weak finite match into a targeted counterexample
and a localized repair obligation---the first index at which the program and the
specification provably differ. Ten of our $984$ also cover entries the encyclopedia marks
\code{dead}: nine are duplicate identifiers, which are redirected to canonical entries,
while one is an erroneous historical version and must not be counted as a current target.
Neither case is evidence that a candidate computes the intended canonical sequence.
Together these
convert the discovery ledger from a count into a graded evidential structure: finite term
agreement, keyword and reference-program corroboration, reviewed formalization, and
finally factive proof or factive counterexample.

\section{From leakage to conjecture}
\label{sec:conjecture}

Near-equivalence is dual to conjecture: what threatens a benchmark as contamination is, read
the other way, a discovered near-identity. The sharpest instance is the pairing of the
\emph{fastest} and the \emph{smallest} candidate covering an entry. These minimize
opposite terms of Levin's $Kt$ complexity,
$Kt(x)=\min_p |p|+\log t(p)$~\cite{levin1973}: program length
and running time. Two unlike programs converging on the same observed terms motivate an
equivalence conjecture, but they provide independent evidence only when their search
lineages are source-disjoint. Otherwise their common training corpus or ancestral
discoveries remain shared causes and additive revision is not licensed. Even with
independent provenance, finite agreement is corroboration rather than entailment.
\emph{Proving} the two programs equivalent on all $n$
converts finite-prefix agreement into total agreement; the induction predicates such proofs
require are exactly what recent conjecturing systems learn to
generate~\cite{gauthier2025}. Our verified layer offers to discharge them as kernel-checked
theorems rather than trusted external calls. The interesting conjectures live where
extensional similarity is high but \emph{syntactic} intensional similarity is low---two
unlike programs apparently computing one sequence---and shared formally proved properties become lemmas
toward the identity proof.

\section{Status and roadmap}
\label{sec:roadmap}

The verified DSL semantics, recognizer, big-step modalities, and the three central obligations
are machine-checked and axiom-clean. The synthesizer's data pipeline, exact split gate, and
the WM-PLN extensional and intensional layers are implemented and measured over the current
corpus, with the numbers reported above; the self-learning loop of Section~\ref{sec:loop}
compounds under three registered replay settings across three worlds. Thirty-three frozen
discoveries now additionally have full-domain correctness proofs against pinned formal
specifications. Near-term work is to (i) independently review the 33 source-to-Lean
translations; (ii) adjudicate a cohort selected from the shortest, weakest-evidence entries,
where counterexamples are most likely; (iii) prove the Python abstract interpreter sound
against the Lean evaluator before using its conclusions for hard pruning; (iv) bind the
Lean canonical policy to exported operator-table metadata; and (v) retain the fastest and
smallest program per sequence for provenance-aware equivalence conjectures. The unifying
aim is an OEIS synthesis loop whose catalogue coverage, source interpretation, and
full-domain program claims are explicitly distinguished and independently auditable.

\begin{thebibliography}{9}
\bibitem{gauthier2023} T.~Gauthier and J.~Urban.
  \emph{Learning Program Synthesis for Integer Sequences from Scratch}.
  AAAI 2023; \href{https://arxiv.org/abs/2202.11908}{arXiv:2202.11908}.
\bibitem{gauthier2025} T.~Gauthier and J.~Urban.
  \emph{Learning Conjecturing from Scratch}.
  \href{https://arxiv.org/abs/2503.01389}{arXiv:2503.01389}, 2025.
\bibitem{oeis} OEIS Foundation Inc.
  \emph{The On-Line Encyclopedia of Integer Sequences} and the \code{oeisdata}
  repository. \url{https://oeis.org}. Data snapshot:
  \url{https://github.com/oeis/oeisdata/tree/a6e0f22854cc1c307da428e9d6295093781df7fa}.
  OEIS content is licensed CC BY-SA 4.0.
\bibitem{cousot1977} P.~Cousot and R.~Cousot.
  \emph{Abstract Interpretation: A Unified Lattice Model for Static Analysis of Programs by
  Construction or Approximation of Fixpoints}. POPL 1977.
  \url{https://doi.org/10.1145/512950.512973}.
\bibitem{levin1973} L.~A.~Levin.
  \emph{Universal Sequential Search Problems}. Problems of Information Transmission,
  9(3):265--266, 1973. \url{https://www.mathnet.ru/eng/ppi914}.
\bibitem{pln2009} B.~Goertzel, M.~Ikl\'e, I.~Goertzel, and A.~Heljakka.
  \emph{Probabilistic Logic Networks}. Springer, 2009.
  \url{https://doi.org/10.1007/978-0-387-76872-4}.
\end{thebibliography}

\end{document}
